\chapter{Requirements Analysis and System Specification}


\section{Stakeholder Analysis and User Research}

\textbf{Primary Stakeholder Profile: Recreational Naturalist:}

\textbf{Demographic and Professional Characteristics:}
\begin{itemize}
    \item \textbf{Representative:} Professional software developer, age 35
    \item \textbf{Technical Proficiency:} Advanced digital literacy
    \item \textbf{Domain Experience:} Intermediate outdoor recreation and nature observation
    \item \textbf{Primary Activities:} Regular ecological excursions, wildlife photography, environmental documentation
    \item \textbf{Research Objectives:} Species identification enhancement, systematic observation logging, knowledge acquisition
    \item \textbf{Technical Challenges:} Taxonomic classification accuracy, data persistence, information management scalability
\end{itemize}

\textbf{Use Case Specification:} Automated Species Recognition and Documentation
\begin{quote}
As a recreational naturalist, the stakeholder requires a computational system capable of automated avian species identification through acoustic pattern recognition, enabling systematic documentation of field observations with integrated photographic evidence and geospatial coordinates within a persistent cloud-based archive.
\end{quote}

\textbf{Functional Requirements Specification:}
\begin{enumerate}
    \item The system shall provide accurate species identification through acoustic signal processing with validated confidence metrics
    \item The system shall support multimedia data capture and cloud synchronization for photographic evidence
    \item The system shall integrate GPS positioning services for precise geospatial documentation
    \item The system shall maintain persistent data storage with cross-platform accessibility through cloud infrastructure
\end{enumerate}

\textbf{Secondary Stakeholder Profile: Advanced Ornithological Practitioner:}

\textbf{Demographic and Professional Characteristics:}
\begin{itemize}
    \item \textbf{Representative:} Retired educational professional, age 50
    \item \textbf{Technical Proficiency:} Moderate digital adoption with learning willingness
    \item \textbf{Domain Experience:} Advanced ornithological knowledge and field experience
    \item \textbf{Primary Activities:} Systematic avian observation, conservation research, peer collaboration
    \item \textbf{Research Objectives:} Digital methodology adoption, comprehensive data management, professional knowledge sharing
    \item \textbf{Technical Challenges:} Traditional-to-digital workflow migration, data accuracy validation, large-scale information management
\end{itemize}

\textbf{Use Case Specification:} Comprehensive Digital Ornithological Platform
\begin{quote}
As an advanced ornithological practitioner, the stakeholder requires a sophisticated digital platform for species identification and systematic documentation, enabling efficient observation recording, comprehensive species data access, and collaborative research data exchange with the ornithological community.
\end{quote}

\textbf{Functional Requirements Specification:}
\begin{enumerate}
    \item The system shall provide professional-grade acoustic species identification with verifiable accuracy metrics
    \item The system shall support comprehensive observation documentation with multimedia evidence and precise geospatial data
    \item The system shall integrate extensive taxonomic databases with detailed species information and behavioral characteristics
    \item The system shall facilitate secure data exchange and collaboration mechanisms with authenticated ornithological practitioners
    \item The system shall maintain comprehensive data archival with long-term accessibility and cross-platform synchronization
\end{enumerate}

\section{Use Case Modeling and Validation}

\textbf{Use Case 1: Archive Access and Data Export Functionality:}
\textbf{Specification:} Multimedia Data Retrieval and External Platform Integration

\begin{quote}
The system shall provide recreational naturalists with intuitive archive navigation capabilities, enabling efficient multimedia content retrieval through advanced search functionality and seamless export integration with external applications, facilitating social media engagement and community knowledge sharing.
\end{quote}

\textbf{Acceptance Criteria and System Requirements:}
\begin{enumerate}
    \item The system shall implement user-friendly archive interface navigation with minimal interaction complexity
    \item The system shall provide comprehensive search functionality with metadata-based filtering for efficient multimedia content discovery
    \item The system shall support cross-platform export capabilities with standardized file format compatibility
    \item The system shall maintain original multimedia quality standards throughout the export process
    \item The system shall ensure persistent cloud-based data storage with reliable accessibility across multiple devices
\end{enumerate}

\textbf{Use Case 2: Real-time Content Sharing and Collaboration:}
\textbf{Specification:} Immediate Multimedia Distribution for Professional Networks

\begin{quote}
The system shall enable advanced ornithological practitioners to rapidly access and distribute recently captured observational evidence through streamlined sharing mechanisms, facilitating immediate collaboration and knowledge dissemination within professional ornithological communities.
\end{quote}

\textbf{Acceptance Criteria and System Requirements:}
\begin{enumerate}
    \item The system shall provide immediate access to recently captured multimedia content through optimized caching mechanisms
    \item The system shall implement integrated sharing functionality directly from content viewing interfaces
    \item The system shall support multi-platform distribution channels including social media, email, and professional communication platforms
    \item The system shall preserve multimedia quality standards throughout the sharing process
    \item The system shall maintain comprehensive data persistence through cloud-based archival systems
\end{enumerate}

\textbf{Use Case 3: Geospatial Data Analysis and Location Intelligence:}
\textbf{Specification:} Species-specific Location Retrieval and Mapping


\section{Functional Requirements}

\begin{enumerate}
    \item \textbf{Audio Classification:} Process bird calls with ≥85\% accuracy in <5 seconds
    \item \textbf{Visual Recognition:} Detect and classify birds from camera images
    \item \textbf{Species Database:} Provide access to 33,000+ bird species with multilingual support
    \item \textbf{Observation Logging:} Record sightings with photos, audio, GPS, and notes
    \item \textbf{Data Management:} Store observations locally with optional cloud sync
    \item \textbf{Offline Operation:} Function completely without internet connectivity
\end{enumerate}

\section{Non-Functional Requirements}

\begin{enumerate}
    \item \textbf{Performance:} App startup <2 seconds, First Start Database under a Minute, UI response <1 second
    \item \textbf{Compatibility:} Android 8.0+ and iOS 12.0+
    \item \textbf{Storage:} Minimum 500MB available space for ML models and Birdex Encyclopedia
    \item \textbf{Privacy:} All ML processing on-device, no data transmission required
\end{enumerate}

